% paper/sections/09c_hfe.tex
%
% §9.4  Hermite 場延長法（HFE）
%   付録D.4から本編§9に昇格 (2026-04-08)
%   理由：分相PPEの基盤技術であり、本編の重要構成要素
%
%   Aslam (2004) の Extension PDE の定常解を CCD Hermite データから
%   直接構成し，界面越し場延長を高次に実現する．
% ═══════════════════════════════════════════════════════════════════════════════
\subsection{\texorpdfstring{Hermite 場延長法（HFE）：$\Ord{h^6}$ 界面越し場延長}{Hermite 場延長法（HFE）：O(h⁶) 界面越し場延長}}
\label{sec:field_extension}
\label{sec:hfe}% canonical label
\label{sec:extension_pde}% backward-compat: §1–§9 cross-refs
\label{sec:gfm}% backward-compat: downstream cross-refs
% ═══════════════════════════════════════════════════════════════════════════════

本節では，\textbf{Hermite 場延長法}（Hermite Field Extension, 以下 \textbf{HFE}）を定式化する．
HFE は Aslam (2004)~\cite{Aslam2004} の Extension PDE の定常解を，
CCD の Hermite データから仮想時間反復を経由せず直接構成する手法である．
1次元または格子軸方向の延長では $(f, f', f'')$ による5次 Hermite 補間で
$\Ord{h^6}$ 精度を得る．2次元の斜め界面では，後述の混合微分データを含む
テンソル積 Hermite 補間を用いると同じ設計精度を維持できる．

\noindent\textbf{ソルバにおける位置づけ．}
低密度比（$\rho_l/\rho_g \leq 5$）の smoothed Heaviside 一括解法では，
圧力場 $p$ は $\varepsilon$-regularized な滑らかな場であり，
HFE は不要であり，適用すると不要な外挿誤差を加える可能性がある
（§~\ref{sec:ppe_discretization_choice}）．
高密度比（$\rho_l/\rho_g > 5$）では\textbf{分相 PPE 解法}
（§~\ref{sec:split_ppe_motivation}；検証は §~\ref{sec:split_ppe_recovery}）を採用し，
界面条件付きの各相 PPE を解く．圧力場が Young--Laplace ジャンプ $[p]_\Gamma = \sigma\kappa$ を
陽に持つため，HFE による界面越し場延長が\textbf{必須}となる．
本節ではこの分相 PPE 解法の\textbf{基盤技術}として HFE を定式化する．

\paragraph{なぜ分相 PPE 解法に界面越し場延長が必要か.}
分相 PPE 解法では各相で定数密度の PPE を独立に解くため，
圧力場が Young--Laplace 則 $[p]_\Gamma = \sigma\kappa$ に従う
不連続を持つ（式~\eqref{eq:young_laplace_jump}）．
CCD 法は同一ステンシル $(i-2,\ldots,i+2)$ で $\Ord{h^6}$ の微分を求解するが，
界面を跨ぐステンシルは Gibbs 振動を引き起こし精度が著しく低下する．

HFE はこの問題を緩和する：
ソース相（例：液相）の場の値を界面法線方向にターゲット相へ延長し，
CCD が界面近傍でも滑らかな場に対して演算できるようにする．
具体的な適用箇所は以下の通りである：
\begin{enumerate}[noitemsep]
  \item \textbf{圧力場の前処理：}
        PPE 求解前に $p^n$ を界面越しに延長し，IPC Predictor の $-\bnabla p^n$ 項から
        Gibbs 振動を除去する．
  \item \textbf{圧力増分の後処理：}
        PPE で得られた $\delta p$ を界面越しに延長し，速度補正
        $\bu^{n+1} = \bu^* - (\Delta t/\rho)\,\bnabla(\delta p_{\mathrm{ext}})$ の精度を維持する．
  \item \textbf{分相 PPE の界面帯処理：}
        界面近傍で片側場を延長し，各相内の CCD 演算が同一相の滑らかな値だけを参照するようにする．
        smoothed Heaviside 一括 PPE の収束性限界については
        §~\ref{sec:varrho_ppe_limitation} で検証する．
\end{enumerate}

\subsubsection{背景：Extension PDE の定常解と HFE の着想}

スカラー場 $q$（圧力増分 $\delta p$ 等）を液相側（$\phi < 0$）から気相側（$\phi \ge 0$）へ滑らかに延長する
偏微分方程式は：
\begin{equation}
  \frac{\partial q}{\partial \tau} + \mathrm{sgn}(\phi)\,\hat{\bm{n}}\cdot\bnabla q = 0
  \label{eq:extension_pde_main}
\end{equation}
ここで $\tau$ は仮想時間（次元は長さ $[L]$；特性速度 $|\hat{\bm{n}}| = 1$ により $\Delta\tau \sim h$），
$\hat{\bm{n}} = \bnabla\phi/|\bnabla\phi|$ は界面法線ベクトルであり，
$\mathrm{sgn}(\phi) = +1$（気相）$/-1$（液相）が延長方向を定める．
液相側（ソース相）の値を固定したまま，気相側（ターゲット相）を仮想時間前進させることで
$\bnabla q \cdot \hat{\bm{n}} = 0$（法線方向勾配ゼロ）の定常解へ収束する．

\subsubsection{CCD との整合性}

最近接点 Hermite 補間（§~\ref{sec:extension_pde_disc}）により構成された延長後の場 $q_{\mathrm{ext}}$ は，
ソース側データと界面ジャンプ条件が同じ次数で与えられる範囲で，
定常解 $\bnabla q \cdot \hat{\bm{n}} = 0$ を高次精度で満たすため，
界面近傍で $C^\infty$ に近い滑らかな分布を持つ（図~\ref{fig:extension_schematic}）．
この滑らかな場 $q_{\mathrm{ext}}$ に対して CCD 演算子 $D^{(1)}, D^{(2)}$ を適用することで：
\begin{itemize}
  \item 界面直近で CCD ステンシルが圧力ジャンプを直接跨ぐことを避けられる
  \item 速度補正 $\bu^{n+1} = \bu^* - (\Delta t/\rho)\,\bnabla(\delta p)$ の離散化誤差が低減される
  \item Balanced-Force 条件（§~\ref{sec:balanced_force}）と整合した圧力勾配評価が実現される
\end{itemize}

\begin{figure}[htbp]
\centering
\begin{tikzpicture}[scale=0.85]
  % 延長前
  \begin{scope}[xshift=-50mm]
    \node[font=\small\bfseries, text=navy] at (0,2.8) {延長前：$\delta p$ に不連続};
    \draw[->] (-2.5,0) -- (2.5,0) node[right,font=\small]{$x$};
    \draw[->] (0,-0.5) -- (0,2.5) node[above,font=\small]{$\delta p$};
    \draw[very thick, blue2, domain=-2.2:0] plot (\x, {1.8});
    \draw[very thick, red2, domain=0:2.2] plot (\x, {0.3});
    \draw[dashed, gray] (0,-0.3) -- (0,2.3);
    \node[font=\scriptsize] at (0,-0.5) {$\Gamma$};
    \node[font=\scriptsize, blue2] at (-1.2, 2.1) {液体};
    \node[font=\scriptsize, red2] at (1.2, 0.6) {気体};
    \draw[thick, orange2, decorate, decoration={zigzag, amplitude=3pt, segment length=5pt}]
      (-0.3, 1.0) -- (0.3, 1.0);
    \node[font=\scriptsize, orange2] at (0, 0.6) {Gibbs};
  \end{scope}
  % 矢印
  \draw[-Stealth, very thick, navy] (-1.5, 1.2) -- (1.0, 1.2)
    node[midway, above, font=\small, navy]{Hermite 延長};
  % 延長後
  \begin{scope}[xshift=50mm]
    \node[font=\small\bfseries, text=navy] at (0,2.8) {延長後：$\delta p_\mathrm{ext}$ は滑らか};
    \draw[->] (-2.5,0) -- (2.5,0) node[right,font=\small]{$x$};
    \draw[->] (0,-0.5) -- (0,2.5) node[above,font=\small]{$\delta p_\mathrm{ext}$};
    \draw[very thick, blue2, domain=-2.2:2.2] plot (\x, {1.8});
    \draw[dashed, gray] (0,-0.3) -- (0,2.3);
    \node[font=\scriptsize] at (0,-0.5) {$\Gamma$};
    \node[font=\scriptsize, blue2] at (1.2, 2.1) {延長済み};
    \node[font=\scriptsize, teal2] at (0, 0.4) {CCD $\bnabla(\delta p_\mathrm{ext})$: $\Ord{h^6}$};
  \end{scope}
\end{tikzpicture}
\caption{最近接点 Hermite 補間による圧力場の界面越し延長（概念図）．
左：CSF 表面張力による Laplace ジャンプが CCD ステンシルを跨ぐと Gibbs 振動が発生．
右：ジャンプ補正済みの片側値を気体側へ延長した後は，
CCD $\bnabla(\delta p)$ が不連続値を直接跨がずに評価される．}
\label{fig:extension_schematic}
\end{figure}

\subsubsection{HFE の離散化：最近接点 Hermite 補間}
\label{sec:extension_pde_disc}

式~\eqref{eq:extension_pde_main} の定常解は最近接点延長
$q_\mathrm{ext}(\bm{x}) = q\bigl(\bm{x}_\Gamma(\bm{x})\bigr)$ に他ならない
（Aslam~(2004)~\cite{Aslam2004}, §2, Proposition 1；$\bnabla q \cdot \hat{\bm{n}} = 0$ の
特性曲線解が最近接点の値の伝播と一致するという主張，
および \cite{Aslam2004} §3 の数値実験で最近接点延長の一致性が確認されている）．
ここで $\bm{x}_\Gamma$ は $\bm{x}$ から界面 $\Gamma$ への最近接点である．
本稿ではこの定常解を，Extension PDE の仮想時間前進を経由せず，
CCD の Hermite データ $(f, f', f'')$ を用いた直接補間により構成する．

\paragraph{最近接点の計算.}
ターゲット側格子点 $\bm{x}$（$\phi(\bm{x}) \geq 0$）に対し，
法線 $\hat{\bm{n}} = \bnabla\phi/|\bnabla\phi|$ を CCD $D^{(1)}$（$\Ord{h^6}$）で評価し，
最近接点を
\begin{equation}
  \bm{x}_\Gamma = \bm{x} - \phi(\bm{x})\,\hat{\bm{n}}(\bm{x})
  \label{eq:closest_point}
\end{equation}
で求める．この式が最近接点を高精度に与えるには，$\phi$ 自身の界面位置誤差，
SDF 誤差 $||\bnabla\phi|-1|$，および $\hat{\bm n}$ の微分誤差が同時に小さい必要がある．
したがって $\bm{x}_\Gamma$ の精度は CCD の $\bnabla\phi$ 精度だけではなく，
CLS 再初期化と界面位置の精度にも律速される．

\paragraph{CCD Hermite 5次補間.}
CCD は各格子点 $x_i$ で $(f_i, f'_i, f''_i)$ を同時に返す．
隣接2点 $x_a, x_b$（$h = x_b - x_a$，いずれもソース側）の
6個のデータから5次 Hermite 多項式を一意に構成する：
\begin{equation}
  P(\xi) = \sum_{k=0}^{5} c_k \,\xi^k, \qquad \xi = \frac{x - x_a}{h} \in [0,1]
  \label{eq:hermite5}
\end{equation}
拘束条件は $P^{(j)}(0) = h^j f_a^{(j)}$，$P^{(j)}(1) = h^j f_b^{(j)}$（$j = 0,1,2$）の6個であり%
\footnote{$h^j$ は $\xi = (x-x_a)/h$ への無次元化に伴うスケール因子：
$P'(\xi) = h f'(x)$，$P''(\xi) = h^2 f''(x)$．}，
係数 $c_0,\ldots,c_5$ は閉形式で得られる．
補間誤差は
\begin{equation}
  |P(x) - f(x)| \leq \frac{\|f^{(6)}\|_\infty}{6!}\,|(x-x_a)^3(x-x_b)^3| = \Ord{h^6}
  \label{eq:hermite5_error}
\end{equation}
であり，CCD 微分精度と同じ $\Ord{h^6}$ を保つ．

\paragraph{片側外挿.}
$\bm{x}_\Gamma$ がソース側最端セル $[x_i, x_{i+1}]$ の外側
（$x_{i+1}$ がターゲット側）に位置する場合，
ソース側のみの2点 $[x_{i-1}, x_i]$ を用いて外挿する．
外挿距離は高々 $h$ であり，式~\eqref{eq:hermite5_error} から
$|(2h)^3 \cdot h^3| = 8h^6$ となるため $\Ord{h^6}$ が維持される．

\paragraph{2次元への拡張.}
テンソル積 Hermite 補間を用いる．
$\bm{x}_\Gamma = (\xi,\eta)$ に対し，各テンソル節点で
\[
  q,\; q_x,\; q_{xx},\; q_y,\; q_{xy},\; q_{xxy},\;
  q_{yy},\; q_{xyy},\; q_{xxyy}
\]
を用意し，$x$ 方向の5次 Hermite 補間を行った後，その補間値・1階微分・2階微分を
$y$ 方向へ再度5次 Hermite 補間する．
これらの混合微分は，まず $x$ 方向 CCD で $q_x,q_{xx}$ を得て，
次に $q,q_x,q_{xx}$ それぞれへ $y$ 方向 CCD を適用する逐次手順で得る．
したがって必要な処理は「1回の CCD 呼び出しで足りる」ものではなく，
各延長帯で軸方向 CCD を複数回適用する．
この完全テンソル版を用いる場合に限り，2次元斜め界面でも Hermite 補間部の
$\Ord{h^6}$ 設計精度を主張できる．混合微分を省略した簡略版は低コストだが，
一般には $\Ord{h^6}$ 保証を持たない．

\begin{tcolorbox}[defbox, title={補足：仮想時間前進による離散化との比較}]
Aslam~(2004)~\cite{Aslam2004} の元論文では，式~\eqref{eq:extension_pde_main} を
1次風上差分＋前進 Euler で仮想時間反復する手法が提案されている．
この手法は構成が単純だが，風上差分の数値拡散により延長帯の精度が $\Ord{h}$ に制限される．
HFE は同一の定常解を CCD データから直接構成することで，
完全テンソル版では Hermite 補間部の $\Ord{h^6}$ 設計精度を保ちつつ，
仮想時間反復を排除する．
\end{tcolorbox}

\subsubsection{HFE の射影ステップへの組み込み}

射影ステップ内では次の2箇所で HFE を適用する：
\begin{enumerate}
  \item \textbf{Step 5a（IPC Predictor の $-\bnabla p^n$ 直前のサブステップ）：}
        前時刻圧力 $p^n$ に HFE を適用して $p^n_{\mathrm{ext}}$ を構築し，
        Predictor の $-\bnabla p^n$ 項に使用する．
  \item \textbf{Step 6a（Corrector 直前；Step 7 前）：}
        圧力増分 $\delta p$ に HFE を適用して $\delta p_{\mathrm{ext}}$ を構築し，
        速度補正 $\bu^{n+1} = \bu^* - (\Delta t/\rho)\,\bnabla(\delta p_{\mathrm{ext}})$ に使用する．
\end{enumerate}
\textbf{配置の要点}：Predictor では既知圧力 $p^n$ を，
Corrector では圧力増分 $\delta p$ をそれぞれ延長し，
界面ジャンプを持つ圧力場を CCD ステンシルが直接跨がないようにする．

\medskip
\noindent\textbf{本節のまとめと位置づけ．}
HFE を本章（§~\ref{sec:ppe_solve}）で定式化する理由は，CCD の Hermite データを基盤とする場延長理論が
PPE 離散化（§~\ref{sec:ccd_poisson_matrix}）および Balanced-Force 条件（§~\ref{sec:balanced_force}）と
密接に関連し，分相 PPE 解法の不可欠な基盤技術であるためである．
§~\ref{sec:split_ppe_motivation} で議論する分相 PPE 解法では，界面上の圧力ジャンプ処理に
HFE が不可欠となり，高密度比シミュレーション（$\rho_l/\rho_g > 10$）で用いる
主要 PPE 戦略として位置づけられる．

% ═══════════════════════════════════════════════════════════════════════════════
